\vspace{0.5em}\noindent\textbf{Chuẩn bị mã nguồn}\par

\vspace{0.5em}\noindent\textbf{Mục tiêu}\par

Chuẩn bị kho lưu trữ mã nguồn (repository) của ứng dụng ở trạng thái sẵn
sàng để phục vụ công tác rà soát kiểm định và thực hiện quy trình triển
khai lặp lại suôn sẻ mà không vô ý sơ sẩy để lọt lộ các thông số bí mật
hay thi ẩu ảnh hưởng lầm sang tài nguyên sản xuất production.

\vspace{0.5em}\noindent\textbf{Bối cảnh kiến
trúc}\par

Kho chứa repository gánh trọn vây kiến trúc cho một mô hình ứng dụng đa
dịch vụ (multi-service application). Hệ thống liên lạc CI/CD sẽ tự chủ
chia nhỏ và gõ lệnh tiến hành cất dựng riêng biệt từng tệp image cho
backend, dịch vụ chat, và (khấu nối nếu bật) dịch vụ AI, tiến ra tung
tản ban bố các bản mô tả kịch bản (Kubernetes manifests), đồng thời cho
thi hành build tĩnh phân chia phần việc cho frontend ra một nhánh ranh
giới độc lập chuyển tản cho kho bộ đôi S3 kết hợp CloudFront.

\vspace{0.5em}\noindent\textbf{Cấu trúc repository}\par

{
\begingroup
\small
\setlength{\tabcolsep}{3pt}
\renewcommand{\arraystretch}{1.15}
\begin{longtable}{@{}
  >{\raggedright\arraybackslash}p{(\linewidth - 2\tabcolsep) * \real{0.5000}}
  >{\raggedright\arraybackslash}p{(\linewidth - 2\tabcolsep) * \real{0.5000}}@{}}
\toprule\noalign{}
\begin{minipage}[b]{\linewidth}\raggedright
Đường dẫn (Path)
\end{minipage} & \begin{minipage}[b]{\linewidth}\raggedright
Mục đích
\end{minipage} \\
\midrule\noalign{}
\endhead
\bottomrule\noalign{}
\endlastfoot
\texttt{backend/\allowbreak{}} & Nơi ngụ của dịch vụ FastAPI, cấu hình dữ liệu
SQLAlchemy models, tệp di cư cơ sở dữ liệu Alembic migrations, mã lệnh
cho cụm worker cùng tập test nghiệm thu \\
\texttt{frontend/\allowbreak{}} & Không gian trang ứng dụng React/Vite, các bộ kết
nối API client, khối thành phần route con cùng dải code test của
frontend \\
\texttt{chat-\allowbreak{}service/\allowbreak{}} & Dịch vụ đàm thoả vận hanh qua Node.js, Express,
Socket.IO, các lớp tiếp giao cơ sở dữ liệu DynamoDB repositories cùng
tệp trung gian Redis adapter \\
\texttt{ai\_\allowbreak{}service/\allowbreak{}} & Cụm dịch vụ đệm adapter chạy trên FastAPI chuyên
phụ trách gọi qua SageMaker cùng luồng mã đọc hiểu và chấm điểm AI \\
\texttt{k8s/\allowbreak{}app/\allowbreak{}} & Hệ bản vẽ manifest gốc tạo lập Kubernetes namespace,
các thông số service, deployments, jobs, dãn dẻo HPA và ngân sách bồi
thăng PDB \\
\texttt{k8s/\allowbreak{}eks/\allowbreak{}} & Hệ manifest đặc thù chuyên gài cất cho môi trường
EKS bao gồm ConfigMap, ServiceAccount, ALB Ingress cùng minh họa mẫu cho
secret \\
\texttt{k8s/\allowbreak{}observability/\allowbreak{}} & Hồ sơ khai báo quy cho luồng giám sát qua
hai tệp kịch bản ServiceMonitor và PrometheusRule \\
\texttt{scripts/\allowbreak{}k8s/\allowbreak{}} & Tập kịch bản tự động hỗ trợ cho thao tác triển
khai cụm tại chỗ hoặc bủa giăng ra EKS \\
\texttt{scripts/\allowbreak{}ci/\allowbreak{}} & Thư viện tệp phụ giúp dây chuyền tác nghiệp trên
GitHub Actions \\
\texttt{scripts/\allowbreak{}aws/\allowbreak{}} & Dải script thao tác trợ lực tài nguyên
CloudFront, ALB, hàng đợi SQS, quản trị cuộn bản phát hành (rollout)
cùng chăm node group \\
\texttt{.\allowbreak{}github/\allowbreak{}workflows/\allowbreak{}cicd.\allowbreak{}yml} & Linh hồn định danh điều hành trọn
dây chuyền thanh tra rà soát và phát triển tiến thẳng cho môi trường
production \\
\texttt{docs/\allowbreak{}architecture/\allowbreak{}} & Thư mục tài liệu quyết định kiến trúc
(ADRs), luật kiểm soát cạnh tranh song song cùng cam đoan hợp đồng thi
hành người tiêu dùng sự kiện \\
\texttt{Dockerfile.\allowbreak{}ai-\allowbreak{}service} & Bản phác vẽ đóng gói container image
cho dịch vụ trung gian AI adapter \\
\texttt{docker-\allowbreak{}compose.\allowbreak{}yml} & Cột mác khởi tung lập trình bủa quanh các
tài nguyên phụ thuộc khi chạy thử tại chỗ \\
\end{longtable}
\endgroup
}

\vspace{0.5em}\noindent\textbf{Nhận mã nguồn dự án được cấp
quyền}\par

Thao tác thực hiện truy nạp vào kho repository chỉ được phép diễn ra
trên cơ quan tài khoản đã sắm sửa ủy quyền hợp pháp:

\begin{Shaded}
\begin{Highlighting}[]
\FunctionTok{git}\NormalTok{ clone https://github.com/Temp{-}orgo/AWS{-}Internship.git}
\BuiltInTok{cd}\NormalTok{ AWS{-}Internship}
\FunctionTok{git}\NormalTok{ status }\AttributeTok{{-}{-}short} \AttributeTok{{-}{-}branch}
\end{Highlighting}
\end{Shaded}

Tôi nghiêm cấm việc rò rỉ bủa lan mọi token bảo mật, chứng chỉ đăng nhập
nhân chứng hoặc những địa chỉ URL từ xa mang gài theo khóa xác thực bên
trong hồ sơ cẩm nang báo cáo này.

\vspace{0.5em}\noindent\textbf{Kiểm tra nhánh và trạng thái
workflow}\par

\begin{Shaded}
\begin{Highlighting}[]
\FunctionTok{git}\NormalTok{ remote }\AttributeTok{{-}v}
\FunctionTok{git}\NormalTok{ status }\AttributeTok{{-}{-}short} \AttributeTok{{-}{-}branch}
\FunctionTok{git}\NormalTok{ log }\AttributeTok{{-}{-}oneline} \AttributeTok{{-}5}
\end{Highlighting}
\end{Shaded}

Kết quả mong đợi:

\begin{itemize}
\tightlist
\item
  Cọc địa chỉ mạng remote kết trói chuẩn xác về kho mã nguồn có bản
  quyền đích xác thuộc dự án.
\item
  Tình trạng cây mã nguồn làm việc (working tree) được thấu đáo 100\%
  ranh giới ngay trước chuỗi búng ra lệnh triển khai.
\item
  Cân mẫn cẩn khôn giữ vẹn toàn, ngăn cản trót lọt việc sửa đè nhầm lên
  các thay đổi nội bồi địa phương không mang dính líu đến dự án.
\end{itemize}

\vspace{0.5em}\noindent\textbf{Xác định các thư viện phụ thuộc của dịch
vụ}\par

\vspace{0.5em}\noindent\textbf{Backend}\par

\begin{Shaded}
\begin{Highlighting}[]
\BuiltInTok{cd}\NormalTok{ backend}
\ExtensionTok{python} \AttributeTok{{-}m}\NormalTok{ pip install }\AttributeTok{{-}r}\NormalTok{ requirements.txt}
\ExtensionTok{python} \AttributeTok{{-}m}\NormalTok{ ruff format }\AttributeTok{{-}{-}check}\NormalTok{ .}
\ExtensionTok{python} \AttributeTok{{-}m}\NormalTok{ ruff check .}
\ExtensionTok{python} \AttributeTok{{-}m}\NormalTok{ mypy app tests}
\ExtensionTok{python} \AttributeTok{{-}m}\NormalTok{ pytest }\AttributeTok{{-}m} \StringTok{"not postgres"}
\end{Highlighting}
\end{Shaded}

Dịch vụ backend tích hợp ranh giới chạy nền Python 3.12 trong dây chuyền
CI, vận dụng bộ khung FastAPI, thư viện SQLAlchemy, trình di cư Alembic,
bộ kiểm tra thuộc tính Pydantic settings, kho móc AWS boto3, cỗ quan
trắc Prometheus client cùng thư viện truy vết OpenTelemetry.

\vspace{0.5em}\noindent\textbf{Dịch vụ chat}\par

\begin{Shaded}
\begin{Highlighting}[]
\BuiltInTok{cd}\NormalTok{ chat{-}service}
\ExtensionTok{npm}\NormalTok{ ci}
\ExtensionTok{npm}\NormalTok{ run check}
\ExtensionTok{npm}\NormalTok{ run lint}
\ExtensionTok{npm}\NormalTok{ run format:check}
\ExtensionTok{npm}\NormalTok{ run test:ci}
\end{Highlighting}
\end{Shaded}

Dịch vụ chat thiết dựng xoay xung quanh Node.js, framework Express, giao
thức thời gian thực Socket.IO, kho lưu DynamoDB, cầu nối Redis adapter,
cùng nhịp tim số liệu giám sát OpenTelemetry và Prometheus metrics.

\vspace{0.5em}\noindent\textbf{Frontend}\par

\begin{Shaded}
\begin{Highlighting}[]
\BuiltInTok{cd}\NormalTok{ frontend}
\ExtensionTok{npm}\NormalTok{ ci}
\ExtensionTok{npm}\NormalTok{ run test}
\ExtensionTok{npm}\NormalTok{ run build}
\end{Highlighting}
\end{Shaded}

Frontend được dệt nên từ quy chuẩn React 18, trình tăng tốc Vite, thư
viện CSS Tailwind CSS, bộ tiếp thu liên lạc Socket.IO client, công cụ
gọi mạng Axios cùng tầng tầng các trạm giao tiếp API client cho từng
tuyến đường riêng.

\vspace{0.5em}\noindent\textbf{Dịch vụ AI}\par

\begin{Shaded}
\begin{Highlighting}[]
\ExtensionTok{docker}\NormalTok{ build }\AttributeTok{{-}f}\NormalTok{ Dockerfile.ai{-}service }\AttributeTok{{-}t}\NormalTok{ internship{-}ai:ci .}
\end{Highlighting}
\end{Shaded}

Bộ trung gian AI adapter tiến hành phơi mở ra bên ngoài các con đường
dẫn \texttt{/\allowbreak{}health/\allowbreak{}ready}, \texttt{/\allowbreak{}parse-\allowbreak{}job}, \texttt{/\allowbreak{}parse-\allowbreak{}cv},
\texttt{/\allowbreak{}match-\allowbreak{}applications}, \texttt{/\allowbreak{}rerank},
\texttt{/\allowbreak{}cv-\allowbreak{}job/\allowbreak{}group-\allowbreak{}score}, kết hợp cùng một chuỗi các luồng chấm điểm
tương đương. Tại môi trường thực tiễn production, nó sẽ gọi sang bộ máy
SageMaker Runtime thực thụ thay vì nặng nề lôi vác nạp trọn cỗ máy model
khổng lồ về ngự ở trong các worker.

\vspace{0.5em}\noindent\textbf{Mẫu tập tin cấu hình môi
trường}\par

Tận dụng những tệp thông tin ví dụ có sẵn để coi như là nguồn tra khảo
tham chiếu chuẩn mực:

{
\begingroup
\small
\setlength{\tabcolsep}{3pt}
\renewcommand{\arraystretch}{1.15}
\begin{longtable}{@{}
  >{\raggedright\arraybackslash}p{(\linewidth - 2\tabcolsep) * \real{0.5000}}
  >{\raggedright\arraybackslash}p{(\linewidth - 2\tabcolsep) * \real{0.5000}}@{}}
\toprule\noalign{}
\begin{minipage}[b]{\linewidth}\raggedright
Tập tin
\end{minipage} & \begin{minipage}[b]{\linewidth}\raggedright
Mục đích
\end{minipage} \\
\midrule\noalign{}
\endhead
\bottomrule\noalign{}
\endlastfoot
\texttt{k8s/\allowbreak{}app/\allowbreak{}secret.\allowbreak{}local.\allowbreak{}example.\allowbreak{}yaml} & Tệp ví dụ minh họa thông số
secret cho cụm Kubernetes tại chỗ \\
\texttt{k8s/\allowbreak{}app/\allowbreak{}secret.\allowbreak{}local.\allowbreak{}yaml.\allowbreak{}example} & Mẫu bản vẽ secret tại chỗ
cho nhà phát triển \\
\texttt{k8s/\allowbreak{}eks/\allowbreak{}secret.\allowbreak{}example.\allowbreak{}yaml} & Hình hài cấu trúc cho các khóa
secret bên trong EKS mà không để lộ nội dung giá trị thật \\
\texttt{backend/\allowbreak{}requirements.\allowbreak{}txt} & Thư viện mã nguồn phụ thuộc phục vụ
runtime backend \\
\texttt{frontend/\allowbreak{}package.\allowbreak{}json} & Hệ script tự động hóa cùng thư viện cho
frontend \\
\texttt{chat-\allowbreak{}service/\allowbreak{}package.\allowbreak{}json} & Cụm script vận hanh và thư viện đi
kèm của dịch vụ chat \\
\end{longtable}
\endgroup
}

Tuyệt đối cấm đưa lên (commit) các tệp biến thực thể \texttt{.\allowbreak{}env},
\texttt{.\allowbreak{}env.\allowbreak{}production}, file cấu hình kubeconfig, chuỗi url kết nối
database, khóa mã bảo an JWT secret, khóa truy nhập AWS, mã thông báo
token GitHub, private key riêng hay các mật mã dịch vụ phát mail SES.

\vspace{0.5em}\noindent\textbf{Lệnh rà soát kiểm định tại
chỗ}\par

Luồng dây chuyền tự động CI trên mạng có thói quen cho thi triển các
chốt trạm tra rà quan trọng như sau:

\begin{Shaded}
\begin{Highlighting}[]
\FunctionTok{bash} \AttributeTok{{-}n} \VariableTok{$(}\FunctionTok{git}\NormalTok{ ls{-}files }\StringTok{\textquotesingle{}*.sh\textquotesingle{}}\VariableTok{)}
\ExtensionTok{python}\NormalTok{ scripts/ci/infrastructure.py}
\ExtensionTok{python} \AttributeTok{{-}m}\NormalTok{ ruff format }\AttributeTok{{-}{-}check}\NormalTok{ backend}
\ExtensionTok{python} \AttributeTok{{-}m}\NormalTok{ ruff check backend}
\ExtensionTok{python} \AttributeTok{{-}m}\NormalTok{ mypy backend/app backend/tests}
\ExtensionTok{python} \AttributeTok{{-}m}\NormalTok{ pytest }\AttributeTok{{-}m} \StringTok{"not postgres"}\NormalTok{ backend}
\ExtensionTok{npm}\NormalTok{ run test }\AttributeTok{{-}{-}prefix}\NormalTok{ frontend}
\ExtensionTok{npm}\NormalTok{ run build }\AttributeTok{{-}{-}prefix}\NormalTok{ frontend}
\ExtensionTok{npm}\NormalTok{ run test:ci }\AttributeTok{{-}{-}prefix}\NormalTok{ chat{-}service}
\ExtensionTok{docker}\NormalTok{ build }\AttributeTok{{-}t}\NormalTok{ internship{-}backend:ci ./backend}
\ExtensionTok{docker}\NormalTok{ build }\AttributeTok{{-}t}\NormalTok{ internship{-}chat:ci ./chat{-}service}
\ExtensionTok{docker}\NormalTok{ build }\AttributeTok{{-}f}\NormalTok{ Dockerfile.ai{-}service }\AttributeTok{{-}t}\NormalTok{ internship{-}ai:ci .}
\FunctionTok{git}\NormalTok{ diff }\AttributeTok{{-}{-}check}
\end{Highlighting}
\end{Shaded}

Trong những kỳ máy trạm lập trình địa phương của bạn thiếu khuyết trình
quản trị Docker, hãy thật thà ghi nhận việc thử nghiệm build đóng ảnh
Docker là chưa chạy rà kiểm định thay vì ngoan cố tự tung tự mác đánh
dấu là đã hoàn thành công xuất sắc.

\vspace{0.5em}\noindent\textbf{Chuẩn bị cho triển
khai}\par

Trước thời khắc khởi búng kích hoạt dây chuyền sản xuất production:

\begin{enumerate}
\def\labelenumi{\arabic{enumi}.}
\tightlist
\item
  Thẩm tra an tâm các biến môi trường và bí mật (secrets) production của
  GitHub đã thọ vị nằm ngắn gũi đầy đủ theo yêu cầu.
\item
  Xác minh vai trò IAM Role của OIDC là
  \texttt{internship-\allowbreak{}github-\allowbreak{}deploy} sẵn lòng cho khoắc nhãn hoán quyền
  thâu vãng bên trong đúng số tài khoản AWS dự án.
\item
  Kịp thời đảm bảo các kho lưu container ECR hiện hữu sẵn, giăng trọn
  chỗ nạp cho backend, chat và cho AI (trong tình huống công tắc triển
  khai AI service thêu bật).
\item
  Xác minh tinh thần trỗi múa trọn vẻ của một hệ sinh thái gồm RDS,
  Redis, DynamoDB, SQS, S3, CloudFront, ALB, Lambda, SES cùng SageMaker.
\item
  Nghiêm rào kiên định giam ngục tham số
  \texttt{PROCESSING\_\allowbreak{}WORKER\_\allowbreak{}ENABLED} ngả ở mức false cho tận mốc ngày
  cỗ máy suy luận AI bám kèm rúng còi thông quan.
\item
  Chọn bấu nấc \texttt{deploy-\allowbreak{}app} để thao diễn đưa mã app lót vào EKS
  và ấn trượt chế độ \texttt{deploy-\allowbreak{}frontend} cho công tác thiêu tung
  giao diện lên cao.
\end{enumerate}

\vspace{0.5em}\noindent\textbf{Các chế độ của quy trình GitHub
Actions}\par

Cấu hình bộ thi kịch bản (workflow) được khảo soát ra lệnh hậu thuẫn cho
chuỗi các cơ cấu tác nghiệp linh động như:

{
\begingroup
\small
\setlength{\tabcolsep}{3pt}
\renewcommand{\arraystretch}{1.15}
\begin{longtable}{@{}
  >{\raggedright\arraybackslash}p{(\linewidth - 2\tabcolsep) * \real{0.5000}}
  >{\raggedright\arraybackslash}p{(\linewidth - 2\tabcolsep) * \real{0.5000}}@{}}
\toprule\noalign{}
\begin{minipage}[b]{\linewidth}\raggedright
Chế độ
\end{minipage} & \begin{minipage}[b]{\linewidth}\raggedright
Mục đích
\end{minipage} \\
\midrule\noalign{}
\endhead
\bottomrule\noalign{}
\endlastfoot
\texttt{validate} & Thực thi kiểm định an nộ không để xô bồ trát dội
sang cụm production \\
\texttt{deploy} & Con đường đưa bản thử ra thực tiễn theo lối mòn truyền
thống nay vẫn được workflow khoan nhượng chấp thuận \\
\texttt{rollout} & Gõ chuông tái khởi động (restart) trọn bộ các
workload trên EKS mà chẳng nhọc sức re-build các ảnh \\
\texttt{restore-\allowbreak{}compute} & Cám ơn gọi đánh đòn chấn hưng dãn sắm lại
trọn dung nấc năng lực cụm máy managed node group \\
\texttt{deploy-\allowbreak{}app} & Dựng nạp ảnh mới đẩy sang kho ECR đồng thời khai
phá tiến vào EKS thuyên chuyển các workload \\
\texttt{deploy-\allowbreak{}public} & Bung trượt bản ghi manifest cắm neo ALB Ingress
ra thế giới bên ngoài \\
\texttt{deploy-\allowbreak{}frontend} & Build kho mã frontend, đồng bộ thư viện lên
kho S3, kêu gọi lệnh xóa bộ đệm (invalidates cache) trên CloudFront \\
\texttt{full} & Tự động thi khát khao nối liền chu trình sản xuất khép
kín một hơi ráo rã trọn bộ từ A sang Z \\
\end{longtable}
\endgroup
}

\vspace{0.5em}\noindent\textbf{Kết quả mong đợi}\par

Kho tàng mã nguồn tựu chung đạt tiêu chí chốt duyệt thành quả khi mà:

\begin{itemize}
\tightlist
\item
  chuỗi tài liệu thư viện bấu vây các dịch vụ thảy đều hoan hỷ tải nạp
  mượt trôi
\item
  trọn các bài kiểm định của backend, frontend cùng dịch vụ chat rềnvang
  đắc thắng băng đọ qua khỏi không gian test quy đúc
\item
  các container image Docker sẵn lòng chịu gõ trúng cấu hình bạt xếp
  thành quả
\item
  những trang manifest mô tả cho Kubernetes xuất xưởng không trầy móng
  xô xát lọt rào châm trích từ cỗ máy kiểm toán script repository
\item
  thông số biến môi trường cùng chuỗi khóa mật secret bên phía GitHub
  ngạo nghễ điểm tên thọ sinh đầy đủ
\item
  luồng ra mắt hệ thống được rành rẽ phân thân rõ mạn: workload nghiệp
  vụ thẳng đường thâm ngục EKS, còn giao diện web lả buông phi thăng ngụ
  kho S3 bấu trói CloudFront
\end{itemize}

\vspace{0.5em}\noindent\textbf{Các lỗi thường
gặp}\par

{
\begingroup
\small
\setlength{\tabcolsep}{3pt}
\renewcommand{\arraystretch}{1.15}
\begin{longtable}{@{}
  >{\raggedright\arraybackslash}p{(\linewidth - 4\tabcolsep) * \real{0.3333}}
  >{\raggedright\arraybackslash}p{(\linewidth - 4\tabcolsep) * \real{0.3333}}
  >{\raggedright\arraybackslash}p{(\linewidth - 4\tabcolsep) * \real{0.3333}}@{}}
\toprule\noalign{}
\begin{minipage}[b]{\linewidth}\raggedright
Triệu chứng
\end{minipage} & \begin{minipage}[b]{\linewidth}\raggedright
Nguyên nhân
\end{minipage} & \begin{minipage}[b]{\linewidth}\raggedright
Hướng khắc phục
\end{minipage} \\
\midrule\noalign{}
\endhead
\bottomrule\noalign{}
\endlastfoot
Lên tiếng ca than \texttt{Cannot\ find\ build-\allowbreak{}and-\allowbreak{}push-\allowbreak{}image.\allowbreak{}sh} & Chẳng
thể dò tìm ra phương mạn cho tập script trợ giúp & Cho lệnh trỏ múa bám
chuẩn đường dẫn \texttt{scripts/\allowbreak{}build-\allowbreak{}and-\allowbreak{}push-\allowbreak{}image.\allowbreak{}sh} hoặc ra quyết
gán cấu hình biến \texttt{BUILD\_\allowbreak{}IMAGE\_\allowbreak{}SCRIPT} \\
Hỏng gục bước tải file tiêu đề Node (Node header) trong tiến trình CI
cho dịch vụ chat & Khâu dịch mã tệp phụ thuộc bản địa (native
dependency) rền nài thèm muốn chuỗi header cài ngay sở tại & Cần mẫn nạp
dải cấu hình biến \texttt{npm\_\allowbreak{}config\_\allowbreak{}nodedir} nhắm trúng vào mạn thư
viện headers mà hệ thống ci toolcache lưu cẩn kiệt \\
Trang frontend build ra bốc sa mù chạy sai về mạn localhost & Mang tiếng
thất đức bỏ quên tham số biến cố định \texttt{VITE\_\allowbreak{}API\_\allowbreak{}BASE\_\allowbreak{}URL=/\allowbreak{}api}
cho luồng production & Nhờng kính nạp trang trọng các thông số biến Vite
ở ngay thời khắc nháy còi chu trình build tĩnh phục vụ production \\
Đóng cửa cự tuyệt lệnh gõ docker build & Tiến trình điều hành Docker
daemon thiêm đố ốm chết im lì giữa máy & Kích ngục hâm ấm gọi còi giục
Docker daemon đứng dậy múa kiếm, hoặc trang nghiêm chốt từ ngữ khai báo
khâu test build chưa qua thực thi \\
Các bí mật (secrets) kinh hãi hiện diện vọt nhảy giữa trang báo rách git
diff & Sơ sót ngập bùn bỏ vây đưa file \texttt{.\allowbreak{}env} hoặc tệp bí mật
sinh ra tạm bở trượt vọt lên mạn staged & Tháo dọn và đá phăng các tệp
bùi bám khóa bảo an xa khỏi mác commit ra lệnh tiễn còi tráo xoay mới
các thông số khóa đã phơi bãi ra xa \\
\end{longtable}
\endgroup
}

\vspace{0.5em}\noindent\textbf{Kết luận}\par

Thư mục chứa đựng mã nguồn ứng dụng qua chặng kiểm nghiệm gay gắt nay
chính thức cán đích sẵn sàng để ra lệnh xuất phát triển khai theo khuôn
nấc quy an nộ bảo mật cao cấp nhất. Mạch phân phối giao diện web buông
rẽ thành chặng bay tiễn tĩnh riêng cho S3 và CloudFront, trong khi khối
workload EKS an nhiên duy trì lộ trình của mình, đồng loạt vây hãm bảo
an cẩn khôn nhờ sức khỏe dây chuyền quản trị thông lượng xác thực GitHub
Actions phối hợp cùng tấm khiên OIDC.
